% ═══ Methods: Encoding and Feature Extraction ═══
\section{Encoding and Feature Extraction}
\label{sec:methods-encoding}

Each MSA is converted to a dense integer matrix using the He~2012 amino-acid
code (values 0--19), with gaps mapped to state~20.  Consensus residues are
computed as the majority non-gap code per column with lowest-code tie-breaking.
Candidate pairs include all $(i,j)$ with $j-i \geq 6$ and both columns having
valid consensus codes.

Pair features comprise: physicochemical group codes ($g_i, g_j \in \{0,\ldots,7\}$),
separation bins ($s \in \{0,1,2,3\}$ for $[6,11)$, $[11,21)$, $[21,31)$,
$[31,\infty)$), consensus residue codes, Gray-Hamming distance, mutual
information, and perplexity ratio.

Two circuit levels are constructed from pooled counts over training proteins:
Level-1 (256 cells) over group$_i \times$ group$_j \times$ separation-bin;
Level-2 (4$\times$1024 cells) over residue-identity pairs padded to
$32\times32$, one map per separation bin.  Cell labelling uses the ternary
rule: ON if contact rate $\geq T \times$ base rate and $n \geq n_{\min}$;
don't-care if $n < n_{\min}$; OFF otherwise, with defaults $T=2$,
$n_{\min}=30$.
